\section{Jism impulsi. Impulsning saqlanish qonuni}

\begin{enumerate}
	
	\item Massasi 4 t bo‘lgan avtobusning tezligi yo‘lning to‘g‘ri chiziqli qismida 5 m/s dan 10 m/s gacha ortdi. Avtobus impulsining o‘zgarishini toping.
	\begin{enumerate}[label=\Alph*)]
		\item 10 kN·s
		\item 20 kN·s
		\item 15 kN·s
		\item 25 kN·s
	\end{enumerate}
	
	\item 300 m/s tezlik bilan gorizontal yo‘nalishda uchayotgan massasi 6 g bo‘lgan o‘q silliq stol ustida yotgan 500 g massali brusokni teshib kiradi va undan 150 m/s tezlikda uchib chiqadi. Brusokning harakat tezligini aniqlang (m/s).
	\begin{enumerate}[label=\Alph*)]
		\item 1,8
		\item 2,2
		\item 0,9
		\item 3,0
	\end{enumerate}
	
	\item Massalari 400 g va 200 g bo‘lgan ikki jism to‘qnashuvdan so‘ng to‘xtab qoldi. Agar birinchi jism tezligi to‘qnashuvdan oldin 0,8 m/s bo‘lgan bo‘lsa, ikkinchi jism tezligi qanday bo‘lgan (m/s)?
	\begin{enumerate}[label=\Alph*)]
		\item 0,8
		\item 2,4
		\item 1,6
		\item 3,2
	\end{enumerate}
	
	\item Massasi 80 g bo‘lgan shar tinch turgan 40 g massali shar tomon harakatlanmoqda. To‘qnashuvni markaziy absolyut noelastik deb hisoblab, ikkinchi shar impulsining o‘zgarishini toping.
	\begin{enumerate}[label=\Alph*)]
		\item 40 mN·s
		\item 65 mN·s
		\item 35 mN·s
		\item 53 mN·s
	\end{enumerate}
	
	\item 2 kg massali bir xil ikkita sharcha bir-biriga qarama-qarshi harakatlanadi. Bir sharchaning tezligi 3 m/s, ikkinchisiniki 7 m/s. Ikkala sharchaning natijaviy impulsining qiymatini toping.
	\begin{enumerate}[label=\Alph*)]
		\item 8
		\item 10
		\item 14
		\item 20
	\end{enumerate}
	
	\item Massasi 1 kg bo‘lgan jism boshlang‘ich tezliksiz 6 $m/s^2$ tezlanish bilan to‘g‘ri chiziqli tekis tezlanuvchan harakat boshladi. Uning 3 m yo‘ldagi impuls o‘zgarishini aniqlang (N·s).
	\begin{enumerate}[label=\Alph*)]
		\item 3
		\item 6
		\item 12
		\item 18
	\end{enumerate}
	
	\item Gorizontal yo‘nalishda 600 m/s tezlik bilan uchayotgan massasi 40 kg bo‘lgan snaryad portlab ikki bo‘lakka: 30 kg va 10 kg bo‘lindi. Katta bo‘lak dastlabki harakat yo‘nalishida 900 m/s tezlikda harakatlandi. Kichik bo‘lakning tezligining modulini toping (m/s).
	\begin{enumerate}[label=\Alph*)]
		\item 150
		\item 450
		\item 300
		\item 600
	\end{enumerate}
	
	\item Massasi 1 kg bo‘lgan jism 1 m radiusli aylana bo‘ylab 2 rad/s burchak tezlik bilan tekis aylanadi. Aylana markazidan jismga o‘tkazilgan radius-vektorning $180^\circ$ ga burilgandagi jism impulsining o‘zgarish modulini toping.
	\begin{enumerate}[label=\Alph*)]
		\item 1
		\item 2
		\item 8
		\item 4
	\end{enumerate}
	
	\item Ikki jism bir-biriga tomon harakatlanmoqda. Biri 3 m/s tezlikda harakatlanmoqda. To‘qnashuvdan so‘ng birinchi jism harakat yo‘nalishida 1,5 m/s tezlikda harakatlanadi. Ular massalari nisbatini toping.
	\begin{enumerate}[label=\Alph*)]
		\item 3
		\item 1,5
		\item 4,5
		\item 2
	\end{enumerate}
	
	\item Massasi 900 kg bo‘lgan zambarakdan 10 kg massali snaryad gorizontga $60^\circ$ burchak ostida 200 m/s tezlikda uchib chiqadi. Otish paytida qurol qanday tezlik oladi (m/s)?
	\begin{enumerate}[label=\Alph*)]
		\item 2,2
		\item 1,1
		\item 0,5
		\item 3,3
	\end{enumerate}
	
	\item Mototsiklchi qiya trekda qayrilishni bajaradi. Agar ishqalanish koeffitsiyenti 0,75, trekning qiyalik burchagi esa gorizontga $45^\circ$ bo‘lsa, harakatning mumkin bo‘lgan maksimal tezligi minimalidan necha marta ortiq?
	\begin{enumerate}[label=\Alph*)]
		\item 3,5
		\item 5
		\item 7
		\item 14
	\end{enumerate}
	
	\item Muzda turgan odam, muz bo‘ylab massasi 0,5 kg bo‘lgan toshni uloqtirdi. Tosh to‘xtaguncha 2 s ichida 20 m yo‘l bosdi. Agar odamning massasi 60 kg bo‘lsa, tosh otish natijasida odam qanday tezlik bilan harakatlana boshlaydi (m/s)?
	\begin{enumerate}[label=\Alph*)]
		\item 0,34
		\item 0,51
		\item 0,08
		\item 0,17
	\end{enumerate}
	
	\item 5 kg massali miltiqdan o‘q uziladi. Merganning yelkasiga tiralmagan miltiqning orqaga tisarilish tezligi mergan miltiqni yelkasiga qattiq tirab turgan holatdagi orqaga tisarilish tezligidan necha marta ortiq bo‘ladi? Merganning massasi 75 kg.
	\begin{enumerate}[label=\Alph*)]
		\item 16
		\item 15
		\item 8
		\item 20
	\end{enumerate}
	
	\item Massasi 200 kg bo‘lgan sol qirg‘oq bo‘ylab 0,8 m/s tezlik bilan suzmoqda. Unga massasi 60 kg bo‘lgan odam qirg‘oqqa tik yo‘nalishda 5 m/s tezlik bilan sakradi. Solning odam bilan birgalikdagi tezligi qanday bo‘lib qoldi (m/s)?
	\begin{enumerate}[label=\Alph*)]
		\item 1,8
		\item 1,3
		\item 2,1
		\item 0,9
	\end{enumerate}
	
	\item Balandligi 1 m bo‘lgan stolning qirrasida massasi 100 g bo‘lgan kichik sharcha turibdi. Unga gorizontal yo‘nalishda 100 m/s tezlik bilan harakatlanayotgan 10 g massali o‘q tegdi va unda tiqilib qoldi. Sharcha stoldan qanday masofada (m) yerga tushadi?
	\begin{enumerate}[label=\Alph*)]
		\item 2
		\item 6
		\item 4
		\item 8
	\end{enumerate}
	
	\item Silliq relslarda turgan aravachaning bir uchidan ikkinchi uchiga relslarga parallel ravishda odam yurib o‘tsa, aravacha yerga nisbatan qanday masofaga siljiydi (m)? Odam massasi 60 kg, aravacha massasi 120 kg, uzunligi 6 m.
	\begin{enumerate}[label=\Alph*)]
		\item 1
		\item 3
		\item 4
		\item 2
	\end{enumerate}
	
	\item Ko‘ndalang kesimi $6 \, cm^2$ bo‘lgan suv oqimi devorga o‘tkazilgan normalga $60^\circ$ burchak ostida urilib, undan elastik qaytmoqda. Suv oqimining tezligi 12 m/s bo‘lsa, devorga ta’sir etuvchi kuchni toping (N).
	\begin{enumerate}[label=\Alph*)]
		\item 86,4
		\item 43,2
		\item 172,8
		\item 64,8
	\end{enumerate}
	
	\item Quroldan 600 m/s tezlik bilan snaryad uchib chiqadi. Agar porox gazlarining o‘rtacha bosim kuchi 2700 kN ga teng bo‘lsa va snaryad stvol ichida 0,002 s harakatlansa, snaryadning massasini aniqlang.
	\begin{enumerate}[label=\Alph*)]
		\item 4,5
		\item 9
		\item 18
		\item 13,5
	\end{enumerate}
	
	\item 1000 kg massali bosqon 1,8 m balandlikdan sandonga tushadi. Zarbning davomiyligi 0,1 s. Urilish noelastik. Sandon va bosqonning o‘zaro ta’sir kuchining (kN da) o‘rtacha qiymatini aniqlang.
	\begin{enumerate}[label=\Alph*)]
		\item 35
		\item 140
		\item 70
		\item 105
	\end{enumerate}
	
	\item 0,5 kg massali koptok vertikalga $60^\circ$ burchak ostida yo‘nalgan 10 m/s tezlik bilan polga urilib sakrab ketadi. Agar urilish 0,1 s davom etsa, koptokning pol bilan o‘zaro ta’sirlashuvining o‘rtacha kuchini toping (N).
	\begin{enumerate}[label=\Alph*)]
		\item 25
		\item 110
		\item 85
		\item 55
	\end{enumerate}
	
	\item 6 km/s tezlik bilan uchayotgan 5 t massali raketa burilish uchun reaktiv dvigatelini yoqadi. Gazlarning tezligi 2 km/s, yoqilg‘i sarfi 10 kg/s. Qanday radiusli (km da) aylana bo‘ylab burilish sodir bo‘ladi?
	\begin{enumerate}[label=\Alph*)]
		\item 1500
		\item 9000
		\item 3000
		\item 4500
	\end{enumerate}
	
	\item Suv oqimi perpendikulyar joylashgan vertikal devorga uriladi va pastga oqib tushadi. Agar oqimning kesim yuzasi $5 \, cm^2$, tezligi esa 8 m/s bo‘lsa, oqim devorga qanday kuch bilan ta’sir qilishini toping (N).
	\begin{enumerate}[label=\Alph*)]
		\item 16
		\item 32
		\item 64
		\item 48
	\end{enumerate}
	
\end{enumerate}



\begin{table}[h!]
	\centering
	\begin{tabular}{|c|c|c|c|c|c|c|c|c|c|c|c|}
		\hline
		№ & J. & № & J. & № & J. & № & J. & № & J. & № & J. \\ \hline
		1 & B & 5 & A & 9 & A & 13 & A & 17 & A & 21 & B \\ \hline
		2 & A & 6 & B & 10 & B & 14 & B & 18 & B & 22 & B \\ \hline
		3 & C & 7 & C & 11 & C & 15 & C & 19 & C & & \\ \hline
		4 & D & 8 & D & 12 & D & 16 & D & 20 & D & & \\ \hline
	\end{tabular}
\end{table}